%------------------------------------------------------------------------------
% \PART 5
%------------------------------------------------------------------------------
\chapter[Аппаратная реализация IP-блока нейронной сети для распознавания рукописных цифр]
{АППАРАТНАЯ РЕАЛИЗАЦИЯ IP-БЛОКА НЕЙРОННОЙ СЕТИ ДЛЯ РАСПОЗНАВАНИЯ РУКОПИСНЫХ ЦИФР}

%------------------------------------------------------------------------------
% \PART 5.2 Text
%------------------------------------------------------------------------------
\section{Описание функциональных блоков}
\hspace*{12.5 mm}Аппаратная реализация нейронной сети включает несколько 
ключевых функциональных блоков:

  1 Блок обработки пикселя.

  2 Блок вычисления тангенса.

  3 Блок вычисления softmax.

  4 Блок памяти изображения.

  5 Блок управляющего автомата.

  6 Блок MAC.

Каждый из этих блоков реализуется в виде аппаратного модуля на HDL.\@

В блоке обработки пикселя осуществляется установка параметра смещения для 
подготовки к вычислению слоя, а также выбирается соответствующий данному этапу 
весовой коэффициент из памяти.

В блоке вычисления тангенса осуществляется аппроксимированный расчет функции 
активации тангенс гиперболический для LST-1 слоя.

Прямая реализация вычисления тангенса будет занимать большое количество 
вычислительных ресурсов. Поэтому для удобства реализации используется 
аппроксимация в соответствии с формулой\cite{TANH}:
\vspace{1mm}
\begin{equation}
    \tanh{(x)}\approx F(x) = 
      \begin{cases}
      \mathrm{sign}(x),         & \mathrm{\text{если}}\;\;  |x|>2, \\
      (1 + \frac{x}{4})\cdot x, & \mathrm{\text{если}}\;\; -2<x<0, \\
      (1 - \frac{x}{4})\cdot x, & \mathrm{\text{если}}\;\;  0<x<2. \\
      \end{cases}
\end{equation}
\vspace{1mm}

В блоке вычисления softmax происходит сравнение входных данных и определение 
наиболее вероятного класса.

В блоке памяти изображения храниться результат обработки на каждом этапе. От 
приемки входных данных, до выхода LST-1 слоя.

Блок управляющего автомата осуществляет управление и синхронизацию между всеми 
блоками устройства и обрабатывает входные управляющие сигналы.

Блок MAC строится на основе блока DSP48. Данный блок используется для 
вычисления операции умножения с накоплением. В данном блоке по сигналу 
инициализации устанавливается смещение, необходимо для текущего этапа 
вычислений, а затем выполнение операции MAC на каждом такте синхронизации.

%------------------------------------------------------------------------------
% \PART 5.4 Text
%------------------------------------------------------------------------------
\section{Описание процесса создания блочной диаграммы}
\hspace*{12.5 mm}Перейдем к построению блочной диаграммы аппаратной части 
проекта в среде Vivado. Блочная диаграмма представляет собой графическое 
описание аппаратных соединений между IP-ядрами, процессорной системой и 
периферией.

Для обеспечения связи между несколькими IP-блоками и процессором в систему 
добавляется модуль AXI Interconnect. Этот модуль играет ключевую роль в 
маршрутизации данных, поступающих от мастер-устройств к одному или нескольким 
ведомым (slave) IP-блокам. Он автоматически адаптирует адресацию, сигналы 
управления и синхронизации, обеспечивая корректный обмен данными в системе.

Также обязательным элементом системы является Processor System Reset. Этот блок 
управляет сигналами сброса всех модулей, входящих в состав блочной диаграммы. 
Он обеспечивает корректную и синхронную инициализацию всех компонентов после 
подачи питания или программного сброса, предотвращая возможные ошибки при 
старте системы.

На рисунке~\ref{fig:bd} показан результирующий вид блочной диаграммы, 
содержащей процессорную систему, пользовательский IP-блок, AXI Interconnect и 
блок сброса. Все соединения между модулями реализованы с соблюдением требований 
интерфейсов AXI и рекомендаций Vivado.

\insertfigure{bd}{pic/bd}{Результирующий вид блочной диаграммы}{16cm}

Для корректной работы системы требуется обеспечить стабильный тактовый сигнал с 
частотой 80 МГц. По умолчанию, тактовый генератор в Zynq-процессорной системе 
может выдавать несколько частот, но для нужд пользовательского IP-блока 
необходима именно частота 80 МГц.

Для этого на вкладке Clock Configuration в настройках Processing System 7 
выполняется установка пользовательской частоты. Один из выходных тактовых 
сигналов (FCLK\_CLK0) настраивается на значение 80 МГц.

%------------------------------------------------------------------------------
% \PART 5.5 Text
%------------------------------------------------------------------------------
\section{Результаты синтеза и симуляции}
\hspace*{12.5 mm}Для проверки работоспособности написанного устройства был 
составлен тест, в котором формируются управляющие воздействия и изображение. 
Работа устройства начинается по сигналу start\_i. 

Для тестирования было взято два изображения, что показано на 
рисунке~\ref{fig:num}.

\insertfigure{num}{pic/num1-4}{Тестируемые изображенияы}{7cm}

Для проверки корректности работы нейронной 
сети необходимо сравнить выход результирующего полносвязного слоя и нейронной 
сети. 

На рисунках~\ref{fig:n1} и~\ref{fig:n4} представлены временные диаграммы 
поведенческого моделирования описанной на SystemVerilog нейронной сети и выход 
эталонной нейронной сети на языке Python.

\insertfigurescustom{pic/t1.jpg}{pic/num1_et.png}{16cm}{Временная диаграмма нейронной сети и результат работы эталона на первом изображении}{n1}{stacked}
\insertfigurescustom{pic/t2.jpg}{pic/num4_et.png}{16cm}{Временная диаграмма нейронной сети и результат работы эталона на втором изображении}{n4}{stacked}

В результате проведённого моделирования и сопоставления работы реализованной 
аппаратной модели с эталонной моделью на языке Python с использованием
библиотеки fixpoint было выявлено полное соответствие выходных данных. Это 
подтверждает корректность проектирования и реализации нейронной сети на 
уровне RTL-описания и её работоспособность при выполнении на целевой ПЛИС.

Для оценки ресурсоёмкости и эффективности реализации проекта был проведён 
анализ использования аппаратных ресурсов, таких как логические элементы 
(LUT), триггеры (FF), специализированные арифметические блоки (DSP) и память 
(BRAM), полученные в процессе синтеза HDL-описания. 

Полученные результаты представлены на рисунке~\ref{fig:hw}.

\insertfigure{hw}{pic/hw1.jpg}{Аппаратные затраты}{11cm}

Для дополнительной верификации корректности функционирования разработанной 
системы проводится временная симуляция на основе результатов синтеза и 
имплементации.

На данном этапе осуществляется моделирование с учетом временных задержек, 
рассчитанных инструментом синтеза, однако ещё без учёта конкретных физических 
ресурсов устройства. 

Постсинтезная симуляция позволяет убедиться, что 
логическая структура проекта сохранена и корректно функционирует после 
преобразования исходного RTL-кода в структуру из примитивов целевой ПЛИС.

Для сохранения структуры блоков при синтезе используются атрибуты, которые 
запрещает среде Xilinx Vivado осуществлять смешивание функциональных блоков
с целью оптимизации:

\small
\fontsize{12pt}{12pt}\selectfont
\begin{verbatim}
  (* dont_touch = "yes" *) (* keep_hierarchy = "yes" *) 
\end{verbatim}
\normalsize
\fontsize{14pt}{14pt}\selectfont

Результаты такой симуляции представлены на рисунке~\ref{fig:post-sin}. Видно, 
что формируемые сигналы соответствуют ожидаемым значениям, соблюдаются 
временные зависимости и интерфейсы взаимодействуют корректно.

\insertfigure{post-sin}{pic/p-sin.jpg}{Временная диаграмма нейронной сети на втором изображении при пост-синтез симмуляции}{16cm}

После завершения этапа имплементации, включающего размещение и трассировку 
(place \& route), производится повторная временная симуляция, уже с учетом всех 
физических задержек, связанных с реальной топологией на кристалле. Это 
позволяет оценить окончательные временные характеристики и убедиться, что 
проект не нарушает временных ограничений (timing constraints).

На рисунке~\ref{fig:post-impl} показана временная диаграмма, соответствующая 
данному этапу. Анализ сигналов подтверждает корректную работу всех блоков 
системы, несмотря на добавление реальных задержек межсоединений и элементов 
логики.

\insertfigure{post-impl}{pic/p-impl.jpg}{Временная диаграмма нейронной сети на первом изображении при пост-имплементация симмуляции}{16cm}

Ещё одним важным параметром при анализе аппаратной реализации является 
потребляемая мощность, которая играет ключевую роль при выборе целевого 
устройства и оценке общей энергоэффективности системы.

На рисунке~\ref{fig:power} представлена оценка энергопотребления 
реализованной нейронной сети на основе результатов имплементации.

\insertfigure{power}{pic/power.jpg}{Оценка потребляемой мощности}{10cm}

Общая мощность составляет 0.166 Вт, из которых:

1 Статическая мощность (Static) — 0.095 Вт, что составляет 57\% от общего 
потребления. Это базовая мощность, потребляемая устройством при включении, вне 
зависимости от его активности.

2 Динамическая мощность (Dynamic) — 0.071 Вт (43\%), обусловлена 
непосредственной работой схем, переключением сигналов, выполнением операций и 
передачей данных.
